\section{Different histories can explain the same closing report}
\label{sec:conditional-history-families}

Knowing a closing balance does not determine the opening balance or the events
that produced it. Consider a deliberately authored monetary example with
closing cash of 15. An opening of 10 followed by receipts of 2 and 3 reaches
that balance. So does an opening of 8 followed by receipts of 3 and 4. The
closing report agrees, but the amounts received during the year are 5 and 7,
and the intermediate balances differ. A simulator that keeps only closing
cash loses information that a later question may require.

The example is a finite representation test. Its two supplied events use an
authored linear template that increases cash, decreases a debit-positive equity
account by the same amount, and increases a monetary cumulative-sales memory
field. The template is not a claim about how an actual issuer should recognize
sales. It makes the distinction between an observed balance and a complete
history explicit while allowing every calculation to be checked by hand.

\subsection{Keep the opening as a shared unknown}

Let \(c_0\) be opening cash at 31 December 2023, and let \(a\) and \(b\)
be the amounts of the supplied events on 30 June and 31 December 2024.
The authored scope gives \(8\leq c_0\leq12\) and
\(0\leq a,b\leq5\). Opening equity is the signed counterpart
\(-c_0\), and cumulative sales start at zero. Under the declared incidence,
the dated states satisfy
\begin{align}
\label{eq:conditional-history-example}
 c_{\mathrm{June}}&=c_0+a, & c_{\mathrm{December}}&=c_0+a+b,\\
 z_{\mathrm{December}}&=a+b, & c_{\mathrm{December}}&=15.
\end{align}
The same \(c_0\) appears at both dates. It cannot be chosen independently
for each observation. Substituting \((c_0,a,b)=(10,2,3)\) gives June cash
12 and cumulative sales 5; substituting \((8,3,4)\) gives June cash 11 and
cumulative sales 7. Both have December cash 15. These are two members of a
continuous conditional family, not two equally likely scenarios or an
enumeration of every possible member.

More generally, let \(x_0\) collect the declared opening account balances
and monetary memory fields. For each supplied event \(j\), let \(u_j\)
be its nonnegative unknown amount, \(t_j\) its date and \(b_j\) its fixed
incidence vector. A snapshot at date \(t\) is
\begin{equation}
\label{eq:conditional-history-affine-state}
 x(t)=x_0+\sum_{j:\,t_j\leq t}b_j u_j.
\end{equation}
The posting components of each \(b_j\) sum to zero, and the opening account
balances also sum to zero. These conditions preserve the balance of the
authored book. The memory components retain the specified monetary movements
without treating them as additional accounts. A snapshot measurement applies
its declared linear form to \(x(t)\) and requires the result to lie within
its supplied interval.

The reference \path{ConditionalMonetaryHistory/1} keeps \(x_0\), event
amounts, domains, templates, dates, measurements and proposed premises in the
family definition. It accepts symbolic opening intervals rather than requiring
one fixed opening assignment. Civil dates must be valid and chronological;
events on the same date retain their supplied order for replay. The finite
additive equations do not establish order-dependent liquidity or recognition
preconditions. Supplied event slots and exclusions define the search scope;
they are not evidence that other transactions did not occur.

Every reported witness is replayed independently with exact fraction arithmetic.
The check starts from the original opening domains, applies each template to
its event amount, verifies balanced postings and account totals, and recomputes
every requested snapshot measurement. An altered after-state cannot substitute
for that replay. The constant-opening fragment also agrees with the preceding
finite reconstruction reference under identical supplied dates, templates,
amounts and measurements.

This first profile uses a single explicit monetary dimension and scale for its
accounts and memory. Inventory units, interest rates, nonlinear valuations and
general contract or cohort structures require richer types and methods. An
affine balanced history can also violate legal or operational conditions that
its supplied constraints do not express. The reference therefore reports no
native execution or accounting admission merely because its equations fit.

\subsection{Attach a measurement family without choosing its witness}

The measurement network from the preceding sections is a family of possible
quantities, not a list of constants waiting to be copied into an opening book.
Let \(\mathcal M\) denote its complete set of assignments \(m\), and let
\(\mathcal H\) denote the authored history family. Explicit mappings
\(g_k\) connect selected measurements to the corresponding history quantities.
The attachment asks for
\begin{equation}
\label{eq:conditional-history-attachment}
 \mathcal J=\left\{(h,m)\in\mathcal H\times\mathcal M:
 m_k=g_k(h)\ \text{for every supplied mapping }k\right\}.
\end{equation}
Both parent formulas remain in this set. The mapping equations add constraints;
they do not replace either parent by a solver-selected assignment. Parent
variables have separate namespaces until a declared equation relates them.
Identical text identifiers in different parents do not create an implicit join.

Suppose the measurement family permits opening cash between 8 and 12, a period
cash change between 3 and 7, closing cash 15, and the opening-plus-change
identity. The assignment \((8,7,15)\) satisfies that family. If a candidate
history instead fixes opening cash at 10, choosing \((8,7,15)\) first would
make the attachment appear impossible. Keeping the full measurement formula
allows the compatible assignment \((10,5,15)\). A feasible parent witness
is therefore insufficient as the input to this join.

The opposite failure is also possible. Keep the history opening fixed at 10,
but constrain the measurement opening to 8. Each parent can separately have a
witness, yet the explicit opening mapping requires \(10=8\). The attached
system is conditionally infeasible. Separate feasibility cannot establish
joint compatibility, just as separate residual intervals could not establish
the earlier measurement network's compatibility.

\subsection{Distinguish a snapshot mapping from a movement mapping}

A snapshot mapping evaluates a declared account or monetary-memory linear form
at the observation's exact instant. A movement mapping instead collects the
signed contributions of events in the observation's half-open period. For
an account component with incidence coefficient \(b_{ij}\), that movement is
\begin{equation}
\label{eq:conditional-history-movement}
 \Delta_i([s,e))=\sum_{j:\,s\leq t_j<e}b_{ij}u_j.
\end{equation}
It excludes the opening balance. In the first authored history, the annual
cash movement is \(2+3=5\), opening cash is 10 and closing cash is 15.
Replacing the movement by the closing snapshot would answer a different
question even though both are monetary quantities.

The mapping must match the observation's value context, dimension, scale and
named quantity. Its snapshot date, or both period endpoints, must be present
in the declared history timeline. A named proposed mapping premise states the
intended population interpretation. Equal amounts or dates do not establish
that interpretation. The checked reference rejects an attempted attachment
of the retained real-source cash-flow context to this synthetic book; no
legal-book cash mapping is inferred from consolidated captions.

The \path{HistoryMeasurementAttachment/1} reference reconstructs both parent
formulas, adds the supplied mappings and solves the resulting conjunction.
Independent verification checks the original temporal measurement equations,
the original history replay and every mapping. Snapshot and period-flow checks
use their respective definitions. Alternative valid joint witnesses are allowed;
their different openings, events and memory remain visible.

\subsection{A candidate witness need not be an admitted history}

Three questions remain separate: does the declared numerical formulation have
a candidate, has its requested search been completed, and are the target's
economic and evidentiary requirements satisfied? Missing proposed premises
leave the numerical question unresolved before solving. An explicit unresolved
source or state obligation can instead permit a witness for the relaxed
formulation while leaving target existence unknown. For example, satisfying
the authored cash equations does not supply an admitted opening or the contract
state needed by a native company method.

An observation without a supplied history mapping stays in the measurement
parent and appears as an unresolved target obligation. It is not removed to
make the attachment complete. Likewise, the history parent's unresolved
obligations and the attachment's additional obligations survive the join.
A complete solver search over the declared linear system cannot discharge
those missing conditions. With all declared mappings and premises supplied,
the result can establish a conditional witness for that limited model while
source admission, accounting admission and native execution remain separate.

The recorded references exercise strict intervals, conflicting measurements,
wrong chronology, unbalanced incidence, unknown premises, unmapped observations,
namespace collisions, altered formulas and forged completeness. Zero computational
budgets return unknown. Infeasibility uses the existing exact solver and its
conflict recheck, which remains distinct from an independent formal proof.
The definitions, complete cases and replay evidence are under
\path{data/accounting-master/}:
\begin{quote}\small
\path{conditional-history-family-20261024}\\
\path{history-measurement-attachment-20261024}
\end{quote}

These steps specify a concrete conditional connection between measurements and
supplied linear histories. Section~\ref{sec:typed-economic-history} extends
the authored history to typed economic quantities and explicit intermediate
constraints. General candidate generation, broader dimensional state, nonlinear
method conditions, legal-book mappings and source or treatment admission still
require their own evidence. The retained issuer's actual history
remains unknown. Its missing contracts, inventory cohorts, tax attributes and
usable liquidity cannot be replaced by either of the authored examples.
